\chapter{Compression}\label{ch:compression}

\begin{center}\itshape
To compress is to predict: every bit saved\\
is a guess about what comes next that turned out right.\\
\upshape(the lesson of this chapter, in our words)
\end{center}

In 1951 Robert Fano, who taught the first course on information theory at MIT, gave his graduate students a choice: they could sit the final examination, or write a term paper on a problem he set. The problem was to find the most efficient binary code for a given list of probabilities. Fano did not tell them that he and Shannon had worked on it without settling it. Their method, now called Shannon--Fano coding, split the list from the top: divide the symbols into two groups of nearly equal probability, give the first group the bit $0$ and the second the bit $1$, and repeat inside each group. It is good, but not always the best.

According to the account he gave to \emph{Scientific American} forty years later, David Huffman worked on the problem for a long time and was about to give up and start revising for the examination when the solution came to him: build the tree from the bottom, by repeatedly joining the two least likely symbols. He later said that had he known that the two founders of the field had struggled with it, he might never have tried. His paper of 1952, four pages long, describes the code that every compressor still contains, from the JPEG format of photographs to the zip files of computers.

Thirty years later two Israeli researchers, Abraham Lempel and Jacob Ziv, found compressors that need no probabilities at all: they learn the statistics of the text while reading it, and on any stationary source they eventually compress as well as if they had known the probabilities in advance. Their algorithms of 1977 and 1978 are behind gzip, the GIF and PNG image formats, and most of the compressed files exchanged on the internet. This chapter proves why all these methods work, and why none of them can compress everything.

\begin{think}
A company advertises a program that shrinks \emph{every} file by at least one byte, without loss. A second program claims only to shrink every file of a million bytes by at least one byte. Which claims are possible? Can a compressed file be compressed again, and again?
\end{think}

\section{Codes that can be decoded}\label{sec:ud}

Tome~3, Chapter~4 worked with \emph{prefix codes}, in which no codeword is a prefix of another. One can decode such a code on the fly. A weaker requirement suffices for decoding in principle. A code $c:A\to\{0,1\}^*$ extends to finite sequences by concatenation, $c(x_1\cdots x_k)=c(x_1)\cdots c(x_k)$; it is \emph{uniquely decodable} if this extension is injective on $A^*$. Every prefix code is uniquely decodable, but not conversely: the code $a\mapsto0$, $b\mapsto01$, $c\mapsto11$ is uniquely decodable (Exercise~\ref{ex:suffix}), although to know whether a $0$ is an $a$ or the start of a $b$ one may have to look arbitrarily far ahead. Do these extra codes allow shorter lengths? Brockway McMillan, a colleague of Shannon at Bell Labs, proved in 1956 that they do not.

\begin{theorem}[McMillan 1956]\label{thm:mcmillan}
The lengths $\ell(x)$ of a uniquely decodable code on a finite alphabet $A$ satisfy Kraft's inequality $\sum_{x\in A}2^{-\ell(x)}\le1$. Consequently (Tome~3, Proposition~4.18) there is a prefix code with the same lengths, and the source coding theorem holds for uniquely decodable codes: their expected length is at least $H(X)$.
\end{theorem}

\begin{proof}
Let $K=\sum_x2^{-\ell(x)}$ and $\ell_{\max}=\max_x\ell(x)$. For every $k\ge1$,
\[
K^k=\sum_{x_1,\dots,x_k\in A}2^{-(\ell(x_1)+\dots+\ell(x_k))}=\sum_{m=1}^{k\ell_{\max}}N_k(m)\,2^{-m},
\]
where $N_k(m)$ is the number of sequences $x_1\cdots x_k$ whose encoding has length $m$. These encodings are distinct binary words of length $m$, by unique decodability, so $N_k(m)\le2^m$ and $K^k\le k\ell_{\max}$. If $K>1$, the left side grows exponentially in $k$ and the right side linearly, which is impossible for large $k$. So $K\le1$. The proof of the source coding theorem (Tome~3, Theorem~4.19(a)) used only Kraft's inequality.
\end{proof}

The trick of raising $K$ to a large power, to turn ``at most $1$'' into ``at most polynomial'', is worth remembering; it appears again in the proof that entropy is additive for Kolmogorov complexity in Chapter~4.

\section{Huffman's code}\label{sec:huffman}

Let $X$ take values in $A=\{x_1,\dots,x_m\}$, $m\ge2$, with probabilities $p_1\ge p_2\ge\dots\ge p_m>0$. A prefix code is the same thing as a binary tree whose leaves are labelled by the letters: the codeword of a letter spells the path from the root to its leaf, $0$ for left and $1$ for right. Its expected length $L=\sum_ip_i\ell_i$ is the average depth of the leaves. Huffman's algorithm: \emph{if $m=2$, give the two letters the codewords $0$ and $1$; otherwise, replace the two least likely letters $x_{m-1}$, $x_m$ by one new letter $z$ of probability $p_{m-1}+p_m$, find a Huffman code for the $m-1$ letters obtained, and give $x_{m-1}$ and $x_m$ the codewords $c(z)0$ and $c(z)1$.}

\begin{figure}[ht]
\centering
\fig{ch02-huffman}
\caption{Huffman's algorithm on the probabilities $0.4$, $0.2$, $0.2$, $0.1$, $0.1$. The two smallest are merged first ($0.1+0.1$), then the two smallest of what remains, and so on; the codewords are read from the root. The expected length is $2.2$ bits, against an entropy of $2.122$ bits.}
\label{fig:huffman}
\end{figure}

\begin{lemma}\label{lem:huffswap}
There is an optimal prefix code (one of least expected length) in which $\ell_1\le\ell_2\le\dots\le\ell_m$, and the codewords of $x_{m-1}$ and $x_m$ have the maximal length and differ only in their last bit.
\end{lemma}

\begin{proof}
Optimal codes exist: a tree with an internal node having a single child can be improved by contracting that edge, and there are finitely many trees with $m$ leaves in which every internal node has two children. Take an optimal code. If $p_i>p_j$ but $\ell_i>\ell_j$, exchanging the codewords of $x_i$ and $x_j$ changes the expected length by $(p_i-p_j)(\ell_j-\ell_i)<0$, which is impossible; and among letters of equal probability we may reorder codewords freely. So we may assume $\ell_1\le\dots\le\ell_m$. In an optimal tree every internal node has two children: otherwise, contract the edge to the only child, which shortens all the codewords below it. So the deepest leaf $x_m$ has a sibling, which is a leaf too (it cannot have descendants deeper than $x_m$), of the same maximal length; exchange the letter on this sibling with $x_{m-1}$, which does not change any length since $\ell_{m-1}=\ell_m$ already.
\end{proof}

\begin{theorem}[Huffman 1952]\label{thm:huffman}
Huffman's code is an optimal prefix code, hence optimal among all uniquely decodable codes.
\end{theorem}

\begin{proof}
Induction on $m$, the case $m=2$ being clear. Let $p'$ be the law on $m-1$ letters obtained by merging $x_{m-1}$ and $x_m$ into $z$. A prefix code $c'$ for $p'$ gives a code $c$ for $p$ by splitting the leaf $z$, and conversely a code $c$ in which $x_{m-1}$ and $x_m$ are siblings gives a code $c'$ by merging them; in both directions
\[
L(c)=L(c')+p_{m-1}+p_m,
\]
since only the two leaves below $z$ are one level deeper. By the induction hypothesis the Huffman code $c'_H$ for $p'$ is optimal, and Huffman's code for $p$ is its split. Let $c^*$ be an optimal code for $p$ as in Lemma~\ref{lem:huffswap}, and $c'$ its merge. Then
\[
L(c_H)=L(c'_H)+p_{m-1}+p_m\le L(c')+p_{m-1}+p_m=L(c^*).\qedhere
\]
\end{proof}

Optimal, but not equal to $H$: on a coin with $p=0.1$, Huffman must spend one bit per toss, while $h(0.1)\approx0.469$. The cure is to code blocks of $n$ symbols, whose Huffman code costs at most $H+\frac1n$ per symbol by Tome~3, Theorem~4.19(c), at the price of an alphabet of size $|A|^n$. For the five probabilities $0.35$, $0.17$, $0.17$, $0.16$, $0.15$, the top-down Shannon--Fano method gives lengths $2,2,2,3,3$ and an expected length of $2.31$ bits, Huffman's $1,3,3,3,3$ and $2.30$ bits: Huffman knew he had beaten his teacher.

\section{Arithmetic coding}\label{sec:arithmetic}

There is a way to code a long sequence without building a code for all sequences: represent the sequence by a number. Order $A=\{1,\dots,m\}$ and give each sequence $x^n=x_1\cdots x_n\in A^n$ the interval
\[
I(x^n)=\bigl[F(x^n),\,F(x^n)+p(x^n)\bigr[\ \subseteq[0,1[,
\]
where $p(x^n)=p(x_1)\cdots p(x_n)$ and $F(x^n)$ is the total probability of the sequences of length $n$ that come before $x^n$ in lexicographic order. The intervals of the $|A|^n$ sequences tile $[0,1[$ in order, and each $I(x^n)$ is cut into the intervals $I(x^nx_{n+1})$ in proportions $p(1),\dots,p(m)$ (Figure~\ref{fig:arith}). So the interval can be computed letter by letter: start from $[0,1[$, and at each letter keep the sub-interval of that letter. The idea goes back to Peter Elias, around 1960; it became practical with the finite-precision versions of Jorma Rissanen and Richard Pasco in 1976.

\begin{figure}[ht]
\centering
\fig{ch02-arith}
\caption{Arithmetic coding of the word $baa$ for the probabilities $p(a)=0.6$, $p(b)=0.3$, $p(c)=0.1$. Each letter keeps its share of the current interval: $[0.6,0.9[$, then $[0.6,0.78[$, then $[0.6,0.708[$, of length $p(baa)=0.108$. Theorem~\ref{thm:arith} takes $\ell=\lceil\log_2\frac1{0.108}\rceil+1=5$ and the dyadic interval $[0.625,0.65625[$, of binary expansion $0.10100$, which lies inside: the codeword is $10100$. (The shorter dyadic interval $[0.625,0.6875[$ of $1010$ also fits; practical coders look for the shortest one.)}
\label{fig:arith}
\end{figure}

\begin{proposition}\label{prop:dyadic}
Every interval $[a,a+L[\,\subseteq[0,1[$ with $L>0$ contains a dyadic interval $[k2^{-\ell},(k+1)2^{-\ell}[$ with $\ell=\lceil\log_2\frac1L\rceil+1$.
\end{proposition}

\begin{proof}
Since $\ell\ge\log_2\frac1L+1$, we have $2^{-\ell}\le\frac L2$. Let $k=\lceil a2^\ell\rceil$. Then $k2^{-\ell}\ge a$ and $k2^{-\ell}<a+2^{-\ell}$, so $(k+1)2^{-\ell}<a+2\cdot2^{-\ell}\le a+L$.
\end{proof}

\begin{theorem}[arithmetic coding]\label{thm:arith}
Map each $x^n\in A^n$ to the $\ell$ binary digits $w$ of $k$ in Proposition~\ref{prop:dyadic} applied to $I(x^n)$, where $\ell=\lceil\log_2\frac1{p(x^n)}\rceil+1$. This is a prefix code on $A^n$ with
\[
\ell(x^n)<\log_2\frac1{p(x^n)}+2,
\]
so for independent symbols its expected length is less than $nH(X)+2$.
\end{theorem}

\begin{proof}
The dyadic interval $J(x^n)$ of $w$ lies inside $I(x^n)$, and these are disjoint for distinct $x^n$. As in the proof of Kraft's inequality (Tome~3, Proposition~4.18), if $w$ were a prefix of $w'$ the interval of $w'$ would lie inside that of $w$; disjointness excludes this. The bound on $\ell$ is clear, and $\mathbb E\log_2\frac1{p(X^n)}=H(X_1,\dots,X_n)=nH(X)$.
\end{proof}

Two features make arithmetic coding the method of choice. It costs at most two bits more than $\log_2\frac1{p(x^n)}$ for the whole sequence, not one bit per block; and nothing in it requires the symbols to be independent or the law to be known in advance. If the probabilities of $x_i$ are given by a \emph{model}, any rule $q(x_i\mid x_1,\dots,x_{i-1})$ that predicts the next symbol from the past, the same construction applied to $q(x^n)=\prod_iq(x_i\mid x^{i-1})$ gives a prefix code of length less than $\log_2\frac1{q(x^n)}+2$. Compression is reduced to prediction: a good model wastes few bits. The finite-precision versions renormalise the interval as soon as its leading bits are decided, emit those bits, and lose a negligible fraction of a bit per symbol.

Here is a model that knows nothing. For binary sequences, Laplace's rule of succession predicts that the next bit is $1$ with probability $\frac{k+1}{i+2}$ after seeing $k$ ones among $i$ bits.

\begin{proposition}[universal coding of coins]\label{prop:laplace}
For Laplace's rule, $q(x^n)=\dfrac1{(n+1)\binom nk}$ for every binary word $x^n$ with $k$ ones. Arithmetic coding with this model has length
\[
\ell(x^n)<nh\Bigl(\frac kn\Bigr)+\log_2(n+1)+2,
\]
and for independent bits with any $P(X_i=1)=\theta$ its expected length is less than $nh(\theta)+\log_2(n+1)+2$.
\end{proposition}

\begin{proof}
When the $i$-th bit is predicted, $i-1$ bits have been seen; if $k_{i-1}$ of them are ones, the bit gets probability $\frac{k_{i-1}+1}{i+1}$ if it is a one and $\frac{i-1-k_{i-1}+1}{i+1}$ if it is a zero. The denominators multiply to $(n+1)!$. The numerators of the ones are $1,2,\dots,k$ in order, and those of the zeros $1,2,\dots,n-k$, whatever the order of the bits. So $q(x^n)=\frac{k!\,(n-k)!}{(n+1)!}$. By Proposition~\ref{prop:binomial}, $\log_2\frac1{q(x^n)}=\log_2(n+1)+\log_2\binom nk\le\log_2(n+1)+nh(\frac kn)$. For the expectation, $h$ is concave, so $\mathbb E\,h(\frac Kn)\le h(\mathbb E\frac Kn)=h(\theta)$ by Jensen.
\end{proof}

Without knowing $\theta$, the coder pays only $\log_2(n+1)+2$ bits more than the entropy of the whole sequence: about $\frac{\log_2n}n$ bits per symbol. This is the price of learning $\theta$; Jorma Rissanen showed in 1984 that $\frac12\log_2n$ per unknown parameter is unavoidable for most $\theta$ (quoted), and the Krichevsky--Trofimov rule $\frac{k+1/2}{i+1}$ achieves it.

\section{Coding the integers}\label{sec:elias}

How should one write down a positive integer $n$ of unknown size, say the length of a file, so that the reader knows where it ends? Writing $n$ in binary takes $\lfloor\log_2n\rfloor+1$ bits, but is not a prefix code. Elias's solution (1975) is to announce the length first.

\begin{proposition}[Elias codes]\label{prop:elias}
Let $b(n)$ be the binary expansion of $n\ge1$ and $|b(n)|=\lfloor\log_2n\rfloor+1$ its length.
\begin{enumerate}[label=(\alph*)]
\item The \emph{$\gamma$ code} $\gamma(n)=0^{|b(n)|-1}b(n)$ is a prefix code of length $2\lfloor\log_2n\rfloor+1$.
\item The \emph{$\delta$ code} $\delta(n)=\gamma(|b(n)|)\,b'(n)$, where $b'(n)$ is $b(n)$ without its leading $1$, is a prefix code of length $\log_2n+2\log_2\log_2n+O(1)$.
\end{enumerate}
\end{proposition}

\begin{proof}
(a) A reader counts the leading zeros, say $j$, then reads the next $j+1$ bits, which begin with $1$ and form $b(n)$; so the end of $\gamma(n)$ is known from its beginning, which is the prefix property. (b) The reader decodes $\gamma(|b(n)|)$, which is self-delimiting by (a), learns the number $|b(n)|-1$ of bits that follow, and puts back the leading $1$. The length is $2\lfloor\log_2|b(n)|\rfloor+1+|b(n)|-1\le\log_2n+2\log_2(\log_2n+1)+1$.
\end{proof}

Since $\sum_n2^{-|\delta(n)|}\le1$ by Kraft, the lengths $\log_2n+2\log_2\log_2n$ are, up to a constant, a convergent ``code length'' for the integers; one cannot use $\log_2n+\log_2\log_2n$ (Exercise~\ref{ex:intcode}), and the true limit, Rissanen's $\log_2^*n=\log_2n+\log_2\log_2n+\log_2\log_2\log_2n+\cdots$, reappears in Chapter~4.

\section{Lempel and Ziv}\label{sec:lz}

Arithmetic coding needs a model. Lempel and Ziv's 1978 algorithm, LZ78, needs none: it builds a dictionary of the phrases it has seen. Parse the binary sequence $x^n$ greedily into phrases $y_1y_2\cdots y_c$, each phrase being the shortest string that has not occurred as an earlier phrase; so every phrase is an earlier phrase (or the empty word) followed by one bit, and all phrases are distinct, except possibly the last. For example,
\[
1011010100010\ \longmapsto\ 1,0,11,01,010,00,10.
\]
Encode each phrase by the index of its prefix among the earlier phrases ($0$ for the empty word) and its last bit: $(0,1)$, $(0,0)$, $(1,1)$, $(2,1)$, $(4,0)$, $(2,0)$, $(1,0)$. With $c$ phrases each index takes $\lceil\log_2c\rceil$ bits, so the code has length
\[
\ell_{\mathrm{LZ}}(x^n)\le c\bigl(\log_2c+2\bigr),
\]
plus the $O(\log n)$ bits needed to announce $n$ (Elias). The decoder rebuilds the dictionary as it goes. The whole question is how small $c$ is.

\begin{lemma}[counting phrases]\label{lem:phrases}
There are numbers $\varepsilon_n\to0$ such that every binary word of length $n$ is parsed by LZ78 into $c\le\frac n{(1-\varepsilon_n)\log_2n}$ phrases.
\end{lemma}

\begin{proof}
At most $2^j$ phrases have length $j$, and at most one phrase repeats, so at most $2^m$ of the $c$ phrases have length less than $m$ (there are $2^m-1$ nonempty words of length less than $m$). The others have length at least $m$, so $n\ge m(c-2^m)$. Assume $c\ge16$ and take $m=\lfloor\log_2c-\log_2\log_2c\rfloor\ge1$, so that $2^m\le\frac c{\log_2c}$ and $m>\log_2c-\log_2\log_2c-1$:
\[
n\ge c\,\log_2c\,(1-\delta_c),\qquad\delta_c=\frac{\log_2\log_2c+1}{\log_2c}+\frac1{\log_2c},
\]
and $\delta_c\to0$. If $c\le\frac n{\log_2^2n}$ there is nothing to prove (for large $n$). Otherwise $\log_2c\ge\log_2n-2\log_2\log_2n$, and $n\ge c(\log_2n-2\log_2\log_2n)(1-\delta_c)$, with $c\to\infty$ as $n\to\infty$; this is the claim with $1-\varepsilon_n=\bigl(1-\frac{2\log_2\log_2n}{\log_2n}\bigr)\bigl(1-\max_{c\ge n/\log_2^2n}\delta_c\bigr)$.
\end{proof}

So $c=O(\frac n{\log n})$ always: even the worst sequence is parsed into long phrases. On a predictable sequence the phrases grow faster still. The next inequality, due to Ziv, ties the number of phrases to the probability of the sequence.

\begin{lemma}[Ziv's inequality]\label{lem:ziv}
Let $p$ be the law of independent bits with $P(X_i=1)=\theta$. If $x^n$ is parsed into $c$ distinct phrases (the last one possibly repeated, which we ignore at the cost of one phrase), then
\[
\log_2\frac1{p(x^n)}\ge c\log_2c-c\log_2\Bigl(\frac nc+1\Bigr)-c\log_2e.
\]
\end{lemma}

\begin{proof}
By independence $p(x^n)=\prod_jp(y_j)$, where $p(y)$ is the probability of the word $y$ in $|y|$ tosses. Let $c_l$ be the number of phrases of length $l$. The phrases of length $l$ are distinct words, so their probabilities add up to at most $1$, and by the inequality of arithmetic and geometric means their product is at most $c_l^{-c_l}$. Hence $\log_2\frac1{p(x^n)}\ge\sum_lc_l\log_2c_l$. Let $\pi_l=\frac{c_l}c$, a law on the lengths with mean $\mu=\frac nc$. Then
\[
\begin{aligned}\sum_lc_l\log_2c_l&=c\log_2c-c\,H(\pi)\eqbr\ge c\log_2c-c\bigl((\mu+1)\log_2(\mu+1)-\mu\log_2\mu\bigr)\end{aligned}
\]
by Exercise~\ref{ex:geometric} of Chapter~1 (the geometric law maximises entropy for a given mean). Finally $(\mu+1)\log_2(\mu+1)-\mu\log_2\mu=\log_2(\mu+1)+\mu\log_2(1+\frac1\mu)\le\log_2(\mu+1)+\log_2e$.
\end{proof}

\begin{theorem}[Ziv and Lempel 1978]\label{thm:lz}
For every $\theta$, if $X_1,X_2,\dots$ are independent bits with $P(X_i=1)=\theta$, then
\[
\frac1n\,\mathbb E\,\ell_{\mathrm{LZ}}(X^n)\longrightarrow h(\theta).
\]
In fact $\ell_{\mathrm{LZ}}(x^n)\le\log_2\frac1{p(x^n)}+n\,r_n$ for every word $x^n$, where $r_n\to0$ does not depend on $x^n$ or $\theta$.
\end{theorem}

\begin{proof}
Combine $\ell_{\mathrm{LZ}}\le c(\log_2c+2)+O(\log n)$ with Ziv's inequality (the repeated last phrase changes $c$ by one, absorbed in the error):
\[
\ell_{\mathrm{LZ}}(x^n)\le\log_2\frac1{p(x^n)}+c\log_2\Bigl(\frac nc+1\Bigr)+c(2+\log_2e)+O(\log n).
\]
By Lemma~\ref{lem:phrases}, $\frac cn\le\frac1{(1-\varepsilon_n)\log_2n}\to0$. The function $t\mapsto t\log_2(\frac1t+1)$ is increasing on $]0,1]$, so $\frac cn\log_2(\frac nc+1)\le\frac{\log_2((1-\varepsilon_n)\log_2n+1)}{(1-\varepsilon_n)\log_2n}\to0$. This proves the second statement. Taking expectations, $\mathbb E\ell_{\mathrm{LZ}}(X^n)\le nh(\theta)+nr_n$; and $\mathbb E\ell_{\mathrm{LZ}}\ge nh(\theta)$ by McMillan's theorem, since LZ78 with the length announced is uniquely decodable.
\end{proof}

The same proof works for every stationary ergodic source, with Ziv's inequality for Markov approximations of the source (Cover and Thomas, Section~13.5); LZ78 is \emph{universal}. The redundancy is of order $\frac{\log\log n}{\log n}$ bits per symbol, much worse than the $\frac{\log n}n$ of Laplace's coder, which however works only for the model it was designed for. Practical compressors combine both ideas: gzip parses with the 1977 variant, LZ77, which points back into a sliding window of recent text, and then codes the pointers with Huffman codes.

\section{What cannot be compressed}\label{sec:incompressible}

Now the box. A lossless compressor is an injective map $C:\{0,1\}^*\to\{0,1\}^*$; the decompressor is its inverse on the image.

\begin{proposition}\label{prop:nocompress}
\begin{enumerate}[label=(\alph*)]
\item No injective map sends every word of length $n$ to a shorter word, and no lossless compressor shortens every file.
\item For every injective $C$ and every $k\ge1$, fewer than a fraction $2^{-k+1}$ of the words of length $n$ are shortened by at least $k$ bits.
\item If $C$ is a prefix code (or uniquely decodable) on $\{0,1\}^n$, then for a uniformly random word the expected length of $C(X)$ is at least $n$.
\end{enumerate}
\end{proposition}

\begin{proof}
(a) There are $2^n$ words of length $n$ and $1+2+\dots+2^{n-1}=2^n-1$ shorter words. If every file were shortened, apply this to the files of length $n$. (b) The words of length at most $n-k$ are $2^{n-k+1}-1<2^{n-k+1}$, so at most that many words of length $n$ can be mapped to them. (c) By McMillan and the source coding theorem, the expected length is at least $H(X)=n$.
\end{proof}

So the first program of the box is impossible, and the second too, by (a). A good compressor shortens the files people actually make, which are a tiny, highly structured minority of all files, and lengthens the rest slightly. Compressing a compressed file again does not help, because a good compressor's output looks random: if it had visible structure, a better model would have used it. Chapter~4 makes ``looks random'' precise for individual strings.

\begin{beauty}
Every lossless compressor is a bet on what the world looks like, and the source coding theorem is the bookmaker's rule: the codeword lengths define a probability law, $2^{-\ell(x)}$, and the expected cost of the bet exceeds the entropy by exactly the divergence between the world and the bet. Huffman found the best bet when the odds are known, arithmetic coding turns any forecaster into a coder at a cost of two bits, and Lempel and Ziv found a gambler who, knowing nothing, learns the odds as fast as they can be learnt.
\end{beauty}

\section{Exercises}

\begin{exercise}[$\star$]\label{ex:twohuffman}
For the probabilities $0.4$, $0.2$, $0.2$, $0.1$, $0.1$, show that breaking ties differently in Huffman's algorithm gives two codes with lengths $2,2,2,3,3$ and $1,2,3,4,4$. Compare their expected lengths and the variances of their lengths.
\end{exercise}

\begin{exercise}[$\star$]\label{ex:suffix}
Show that the code $a\mapsto0$, $b\mapsto01$, $c\mapsto11$ is uniquely decodable but not a prefix code. How far ahead must a decoder look to decode the first letter of $0111\cdots1$?
\end{exercise}

\begin{exercise}[$\star$]\label{ex:fullnode}
Show that in an optimal prefix code every internal node of the tree has two children, so that $\sum_x2^{-\ell(x)}=1$, and that the lengths are at most $m-1$ for $m$ letters.
\end{exercise}

\begin{exercise}[$\star$]\label{ex:dyadichuff}
If all the $p_i$ are powers of $\frac12$, show that Huffman's code has $\ell_i=\log_2\frac1{p_i}$ and expected length exactly $H$.
\end{exercise}

\begin{exercise}[$\star\star$]\label{ex:fibonacci}
Show that the probabilities proportional to the Fibonacci numbers $F_1,F_1,F_2,\dots,F_{m-1}$ (with $F_1=F_2=1$) can lead to a Huffman tree of depth $m-1$, the largest possible.
\end{exercise}

\begin{exercise}[$\star\star$]\label{ex:blocks}
For independent bits with $P(1)=0.1$, compute the expected length per bit of Huffman's code on blocks of $1$, $2$ and $3$ bits, and compare with $h(0.1)\approx0.469$.
\end{exercise}

\begin{exercise}[$\star\star$]\label{ex:intcode}
Show that $\sum_{n\ge2}2^{-(\log_2n+\log_2\log_2n)}=\sum_{n\ge2}\frac1{n\log_2n}$ diverges, so that no prefix code for the integers has lengths at most $\log_2n+\log_2\log_2n$ for all large $n$, while $\sum_n\frac1{n(\log_2n)^2}$ converges.
\end{exercise}

\begin{exercise}[$\star$]\label{ex:lzparse}
Parse $0001101001110000$ with LZ78 and write its encoding as pairs.
\end{exercise}

\begin{exercise}[$\star\star$]\label{ex:twopart}
A \emph{two-part code} for a binary word $x^n$ with $k$ ones first writes $k$ with $\lceil\log_2(n+1)\rceil$ bits, then the rank of $x^n$ among the $\binom nk$ words with $k$ ones. Show that its length is within $2$ bits of the length of the Laplace arithmetic coder of Proposition~\ref{prop:laplace}, and explain why.
\end{exercise}

\begin{exercise}[$\star\star$]\label{ex:mismatch}
A Huffman code is designed for the law $q$ but the symbols follow the law $p$. Show that the expected length of a code with lengths $\lceil\log_2\frac1{q(x)}\rceil$ lies between $H(p)+D(p\,\|\,q)$ and $H(p)+D(p\,\|\,q)+1$. This is the price of a wrong model.
\end{exercise}

\begin{exercise}[$\star\star\star$]\label{ex:sfe}
\emph{Shannon--Fano--Elias.} Let $p>0$ on $A=\{1,\dots,m\}$, $F(x)=\sum_{y<x}p(y)$ and $\bar F(x)=F(x)+\frac{p(x)}2$. Show that the first $\lceil\log_2\frac1{p(x)}\rceil+1$ bits of the binary expansion of $\bar F(x)$ form a prefix code, of expected length less than $H+2$, without sorting the probabilities.
\end{exercise}

\section{Hints}

\begin{hint}{ex:twohuffman}
After merging $0.1+0.1$, three letters have probability $0.2$.
\end{hint}
\begin{hint}{ex:suffix}
Read the words backwards.
\end{hint}
\begin{hint}{ex:fullnode}
Contract an edge to an only child.
\end{hint}
\begin{hint}{ex:dyadichuff}
The two least likely letters have the same probability; induct.
\end{hint}
\begin{hint}{ex:fibonacci}
$F_1+\dots+F_k=F_{k+2}-1$: the merged node is always one of the two smallest.
\end{hint}
\begin{hint}{ex:blocks}
Run Huffman on $2$, $4$ and $8$ probabilities.
\end{hint}
\begin{hint}{ex:intcode}
Compare with $\int\frac{dx}{x\ln x}$.
\end{hint}
\begin{hint}{ex:lzparse}
The first phrases are $0$, $00$, $1$.
\end{hint}
\begin{hint}{ex:twopart}
Compute $\log_2\frac1{q(x^n)}$ for Laplace's rule.
\end{hint}
\begin{hint}{ex:mismatch}
$\ell(x)-\log_2\frac1{p(x)}$ lies between $\log_2\frac{p(x)}{q(x)}$ and that plus $1$.
\end{hint}
\begin{hint}{ex:sfe}
The intervals $[F(x),F(x)+p(x)[$ are disjoint; truncating $\bar F(x)$ to $\ell$ bits stays inside.
\end{hint}

\section{Solutions}

\begin{solution}{ex:twohuffman}
\idea Huffman's algorithm only says ``merge two of the least likely''; when there are ties, every choice is optimal, but the shapes of the trees differ.
\writeup Call the letters $x_1,\dots,x_5$. First merge $x_4$ and $x_5$ into a node $z$ of weight $0.2$; now $x_2$, $x_3$ and $z$ all weigh $0.2$. If we merge $x_2$ and $x_3$, the weights are $0.4$ ($x_1$), $0.4$ ($x_2x_3$) and $0.2$ ($z$); merging $z$ with $x_1$ and then the two nodes of weight $0.6$ and $0.4$ gives the lengths $2,2,2,3,3$ of Figure~\ref{fig:huffman}. If instead we merge $z$ with $x_3$, the weights are $0.4$ ($x_1$), $0.4$ ($x_3z$) and $0.2$ ($x_2$); merging $x_2$ with $x_3z$, then the result with $x_1$, gives lengths $1,2,3,4,4$. (Merging $z$ with $x_2x_3$ in the first branch gives a third optimal code, $1,3,3,3,3$.) All have expected length $2.2$: for instance $0.4\cdot1+0.2\cdot2+0.2\cdot3+0.1\cdot4+0.1\cdot4=2.2$. The variances of the lengths are $5-2.2^2=0.16$ and $6.2-2.2^2=1.36$: the first code is preferable when buffers must be small, which is why careful implementations break ties in favour of older nodes.
\end{solution}

\begin{solution}{ex:suffix}
\idea A code in which no codeword is a \emph{suffix} of another can be decoded from right to left.
\writeup No codeword among $0$, $01$, $11$ is a suffix of another ($0$ is not a suffix of $01$ or $11$, and $01$, $11$ have the same length). So reading an encoded word from the end, the last codeword is determined, remove it and repeat: the decoding is unique. It is not a prefix code, since $0$ is a prefix of $01$. For $01^k$, the first letter is $a$ if $k$ is even (then $0\,11\cdots11$) and $b$ if $k$ is odd ($01\,11\cdots11$): the decoder must read all the ones before deciding, as far ahead as the word goes.
\end{solution}

\begin{solution}{ex:fullnode}
\idea A node with one child wastes a bit on every codeword below it.
\writeup If an internal node has a single child, replace the subtree at the node by the subtree at the child: every codeword below loses one bit, the code stays prefix-free, and the expected length strictly decreases (the subtree contains a leaf of positive probability). So an optimal tree is full, and in a full binary tree $\sum_{\text{leaves}}2^{-\text{depth}}=1$ (induction: replacing a leaf by two children keeps the sum). A full binary tree with $m$ leaves has $m-1$ internal nodes, and a path from the root passes through at most $m-1$ of them, so depths are at most $m-1$.
\end{solution}

\begin{solution}{ex:dyadichuff}
\idea With dyadic probabilities the ideal lengths $\log_2\frac1{p_i}$ are integers and satisfy Kraft with equality; Huffman finds them because merging the two smallest keeps the law dyadic.
\writeup Induction on $m$; if $m=2$, then $p_1=p_2=\frac12$. Let $m>2$ and let $2^{-a}$ be the smallest probability. The numbers $2^ap_i$ are integers summing to $2^a$, an even number, and they are even except those equal to $1$; so the number of letters of probability $2^{-a}$ is even, hence at least two. Huffman merges two of them into a letter of probability $2^{-a+1}$, which gives a dyadic law on $m-1$ letters. By induction the merged letter gets length $a-1$ and every other letter $\log_2\frac1{p_i}$; the two merged letters get length $a=\log_2\frac1{2^{-a}}$. So $\ell_i=\log_2\frac1{p_i}$ for all $i$, and $L=H$.
\end{solution}

\begin{solution}{ex:fibonacci}
\idea Fibonacci numbers grow so slowly that the newly merged node is always one of the two smallest, so the tree is a long path with one leaf hanging at each level.
\writeup Use the weights $F_1,F_1,F_2,\dots,F_{m-1}$ (normalising does not change the algorithm). We show that after $j$ merges, $1\le j\le m-2$, the weights are the merged node, of weight $F_{j+2}$, and the untouched letters $F_{j+1},F_{j+2},\dots,F_{m-1}$. For $j=1$ the two smallest weights $F_1,F_1$ are merged into $2=F_3$, and $F_2,\dots,F_{m-1}$ remain. If the claim holds for $j$, the two smallest weights are $F_{j+1}$ and $F_{j+2}$; when the merged node ties with a letter, we break the tie in favour of the merged node (this is allowed: both choices are optimal). Their sum is $F_{j+3}$, and the remaining letters are $F_{j+2},\dots,F_{m-1}$, as claimed. Each merge puts the merged node one level below the new one, so the two letters merged first end at depth $m-1$, the maximum possible (Exercise~\ref{ex:fullnode}).
\end{solution}

\begin{solution}{ex:blocks}
\idea Huffman on blocks: the waste of up to one bit per block is spread over the block.
\writeup For one bit, the code must spend $1$ bit. For pairs, the probabilities are $0.81,0.09,0.09,0.01$; Huffman gives lengths $1,2,3,3$ and $0.81+0.18+0.27+0.03=1.29$ bits per pair, $0.645$ per bit. For triples, the probabilities are $0.729$, three times $0.081$, three times $0.009$ and $0.001$; Huffman gives an expected length of $1.598$ bits per triple, $0.533$ per bit (notebook). The values decrease towards $h(0.1)\approx0.469$, slowly: blocks of $n$ bits guarantee only $h+\frac1n$.
\end{solution}

\begin{solution}{ex:intcode}
\idea Kraft turns code lengths into a series that must converge; the integral test decides convergence.
\writeup $2^{-(\log_2n+\log_2\log_2n)}=\frac1{n\log_2n}=\frac{\ln2}{n\ln n}$. Since $x\mapsto\frac1{x\ln x}$ is decreasing, $\sum_{n=2}^N\frac1{n\ln n}\ge\int_2^{N+1}\frac{dx}{x\ln x}=\ln\ln(N+1)-\ln\ln2\to\infty$. A prefix code with $\ell(n)\le\log_2n+\log_2\log_2n$ for all $n\ge n_0$ would have $\sum_n2^{-\ell(n)}=\infty$, contradicting Kraft. On the other hand $\frac1{n(\log_2n)^2}=\frac{(\ln2)^2}{n(\ln n)^2}$ and $\int_2^\infty\frac{dx}{x(\ln x)^2}=\frac1{\ln2}<\infty$, so that series converges; the $\delta$ code realises lengths $\log_2n+2\log_2\log_2n+O(1)$.
\end{solution}

\begin{solution}{ex:lzparse}
\idea Each new phrase is the shortest string not yet in the dictionary: extend a known phrase by one bit.
\writeup Reading from the left: $0$ is new; then $0$ is known and $00$ new; then $1$ is new; then $1$ is known and $10$ new; then $1$, $10$ are known and $100$ new; then $1$ is known and $11$ new; then $1$, $10$, $100$ are known and $1000$ new; a last $0$ remains, which repeats the first phrase. So the parse is $0|00|1|10|100|11|1000|0$ (lengths $1+2+1+2+3+2+4+1=16$), with dictionary indices $1$ to $7$, and the encoding is $(0,0)$, $(1,0)$, $(0,1)$, $(3,0)$, $(4,0)$, $(3,1)$, $(5,0)$, followed by the pointer to phrase $1$ for the repeated last phrase.
\end{solution}

\begin{solution}{ex:twopart}
\idea Both codes spend $\log_2(n+1)$ bits on the number of ones and $\log_2\binom nk$ on the rest; Laplace's rule is the two-part code in sequential disguise.
\writeup The two-part code has length $\lceil\log_2(n+1)\rceil+\lceil\log_2\binom nk\rceil$, which is within $2$ bits of $\log_2(n+1)+\log_2\binom nk$. Laplace's arithmetic coder has length within $2$ bits of $\log_2\frac1{q(x^n)}=\log_2(n+1)+\log_2\binom nk$ (proof of Proposition~\ref{prop:laplace}). The reason is that $q$ is the uniform mixture of the laws of independent coins, $q(x^n)=\int_0^1\theta^k(1-\theta)^{n-k}\,d\theta=\frac{k!(n-k)!}{(n+1)!}$: it gives each count $k$ total probability $\frac1{n+1}$ and spreads it uniformly over the $\binom nk$ words with $k$ ones, which is exactly what the two-part code does.
\end{solution}

\begin{solution}{ex:mismatch}
\idea Write the length of each codeword as ideal length for $p$ plus a correction; the average correction is the divergence.
\writeup With $\ell(x)=\lceil\log_2\frac1{q(x)}\rceil$, we have $\log_2\frac1{q(x)}\le\ell(x)<\log_2\frac1{q(x)}+1$. Average under $p$: $\sum_xp(x)\log_2\frac1{q(x)}=\sum_xp(x)\log_2\frac1{p(x)}+\sum_xp(x)\log_2\frac{p(x)}{q(x)}=H(p)+D(p\,\|\,q)$. So the expected length lies in $[H(p)+D(p\,\|\,q),\,H(p)+D(p\,\|\,q)+1[$. Believing the wrong law costs $D(p\,\|\,q)$ bits per symbol, which is why the divergence is also called the \emph{redundancy} of a code.
\end{solution}

\begin{solution}{ex:sfe}
\idea The midpoint of the interval of $x$ identifies $x$, and with enough bits its truncation cannot leave the interval.
\writeup Let $\ell(x)=\lceil\log_2\frac1{p(x)}\rceil+1$, so $2^{-\ell(x)}\le\frac{p(x)}2$, and let $w(x)$ be the first $\ell(x)$ bits of $\bar F(x)$, read as the number $\lfloor\bar F(x)\rfloor_{\ell(x)}$. Then $\bar F(x)-2^{-\ell(x)}<\lfloor\bar F(x)\rfloor_{\ell(x)}\le\bar F(x)$, so the dyadic interval of $w(x)$, which is $\bigl[\lfloor\bar F(x)\rfloor_{\ell(x)},\lfloor\bar F(x)\rfloor_{\ell(x)}+2^{-\ell(x)}\bigr[$, lies inside $\bigl]\bar F(x)-\frac{p(x)}2,\bar F(x)+\frac{p(x)}2\bigr[\subseteq[F(x),F(x)+p(x)[$. These intervals are disjoint for distinct $x$, so, as in Theorem~\ref{thm:arith}, no $w(x)$ is a prefix of another. The expected length is $\sum p(x)\ell(x)<\sum p(x)(\log_2\frac1{p(x)}+2)=H+2$. This is arithmetic coding for a single symbol.
\end{solution}

\section*{Sources and further reading}

\begin{itemize}
\item \emph{A research paper to read.} J.~Ziv and A.~Lempel, \emph{Compression of individual sequences via variable-rate coding}, IEEE Transactions on Information Theory 24 (1978), 530--536. The paper of LZ78. The authors measure the compressibility of an \emph{individual} infinite sequence by the best compression achievable by finite-state machines, and prove that their parsing algorithm reaches it for every sequence. Read Section~II for the model, then the description of the incremental parsing in Section~III, and compare their counting lemma with our Lemma~\ref{lem:phrases}; the probabilistic statement, our Theorem~\ref{thm:lz}, appears as a corollary for stationary sources. Notice that no probability appears in the main theorem: the paper is a bridge between Shannon's theory and the individual sequences of Chapter~4.
\item \emph{The stories.} G.~Stix, \emph{Profile: David A. Huffman}, Scientific American 265 (1991), no.~3, 54, 58, tells the term-paper story in Huffman's words. D.~A.~Huffman, \emph{A method for the construction of minimum-redundancy codes}, Proceedings of the IRE 40 (1952), 1098--1101; R.~M.~Fano, \emph{The transmission of information}, Technical Report 65, MIT Research Laboratory of Electronics, 1949.
\item B.~McMillan, \emph{Two inequalities implied by unique decipherability}, IRE Transactions on Information Theory 2 (1956), 115--116; the short proof by powers is due to J.~Karush, \emph{A simple proof of an inequality of McMillan}, IRE Transactions on Information Theory 7 (1961), 118. P.~Elias, \emph{Universal codeword sets and representations of the integers}, IEEE Transactions on Information Theory 21 (1975), 194--203.
\item J.~Rissanen, \emph{Generalized Kraft inequality and arithmetic coding}, IBM Journal of Research and Development 20 (1976), 198--203; R.~Pasco, \emph{Source coding algorithms for fast data compression}, PhD thesis, Stanford University, 1976; I.~H.~Witten, R.~M.~Neal and J.~G.~Cleary, \emph{Arithmetic coding for data compression}, Communications of the ACM 30 (1987), 520--540, with a complete program. J.~Rissanen, \emph{Universal coding, information, prediction, and estimation}, IEEE Transactions on Information Theory 30 (1984), 629--636; R.~E.~Krichevsky and V.~K.~Trofimov, \emph{The performance of universal encoding}, IEEE Transactions on Information Theory 27 (1981), 199--207.
\item J.~Ziv and A.~Lempel, \emph{A universal algorithm for sequential data compression}, IEEE Transactions on Information Theory 23 (1977), 337--343 (LZ77); T.~A.~Welch, \emph{A technique for high-performance data compression}, Computer 17 (1984), no.~6, 8--19 (LZW, the variant used in GIF); P.~Deutsch, \emph{DEFLATE Compressed Data Format Specification version 1.3}, RFC~1951, 1996 (the format of gzip and zip).
\item T.~M.~Cover and J.~A.~Thomas, \emph{Elements of Information Theory}, 2nd ed., Wiley, 2006, Chapters~5 and~13; D.~J.~C.~MacKay, \emph{Information Theory, Inference, and Learning Algorithms}, Cambridge University Press, 2003, Chapters~5 and~6, freely available from its author, with a very readable treatment of arithmetic coding as prediction.
\end{itemize}
